\documentclass[UTF8,a4paper,11pt]{ctexart}
\usepackage[margin=2cm]{geometry}
\usepackage{amsmath,amssymb,fancyhdr,hyperref}
\hypersetup{colorlinks=true,urlcolor=blue}
\pagestyle{fancy}\fancyhf{}\fancyfoot[C]{\thepage}\renewcommand{\headrulewidth}{0pt}
\setlength{\parindent}{0pt}\setlength{\parskip}{6pt}
\begin{document}
\begin{center}{\LARGE 竞赛选题·教师解析}\par 新高一开学三周拓展\end{center}
本卷共100分，建议90分钟，后两题也可留作课后挑战。评分为教学用分步评分，不采用原赛事评分标准。难度为选题判断，未经学生实测。解法只依赖集合、初中代数、正数基本不等式和平方非负；第4题需要引导发现计数的存在性证明，第5题换元的思路跨度较大。

\textbf{1．答案：$A=\{-3,0,2,6\}$。（16分）}

令 $T=a_1+a_2+a_3+a_4$。四个三元素子集的和分别为 $T-a_i$，它们互不相同，每个 $a_i$ 在四个和中恰好出现三次，故
\[3T=-1+3+5+8=15,\qquad T=5.\]
所以 $A=\{5-(-1),5-3,5-5,5-8\}=\{-3,0,2,6\}$。代回可得四个和依次为 $-1,3,5,8$，符合题意。

评分：表示四个和4分；整体求和并得 $T=5$ 共6分；恢复集合并检验6分。

目标与评价：集合无序性、整体求和；观察能否摆脱逐项配对。易错：随意对应 $a_i$ 与 $B$ 中元素，或把四个和的总和当作 $T$。难度：中等。

\textbf{2．答案：$M=5$，$m=2\sqrt3$，$M-m=5-2\sqrt3$。（18分）}

因为 $a\ge1,b\le2$，所以 $3/a+b\le3+2=5$；当 $(a,b)=(1,2)$ 时取等，且满足全部约束，故 $M=5$。

另一方面，由 $b\ge a>0$，
\[\frac3a+b\ge\frac3a+a\ge2\sqrt3.\]
两步同时取等要求 $b=a$ 且 $a=\sqrt3$。此时 $1\le a=b=\sqrt3\le2$，故 $m=2\sqrt3$，从而 $M-m=5-2\sqrt3$。

评分：最大值的界及取等6分；最小值的界及取等8分；差值4分。

目标与评价：约束转化、基本不等式的等号条件；必须同时验证界与可达性。易错：只得到一个下界就宣布最小值。难度：中等。

\textbf{3．答案：$(a,b)=(\sqrt2-1,\sqrt2+1)$ 或其交换。（20分）}

令 $s=a+b>0$，$p=ab>0$。第一个条件给出 $s\le2\sqrt2p$；第二个条件给出
\[s^2=(a-b)^2+4ab=4p^3+4p\ge8p^2.\]
因此 $s\ge2\sqrt2p$，两边夹住，必有 $s=2\sqrt2p$，且 $p^3+p\ge2p^2$ 取等。由 $p^3=p$ 及 $p>0$，得 $p=1$、$s=2\sqrt2$。于是 $a,b$ 为 $t^2-2\sqrt2t+1=0$ 的两根，即 $\sqrt2\pm1$。两种次序均满足原条件。

评分：引入和积并转化4分；得到相反方向的界6分；锁定等号得和积4分；求出两对解并检验6分。

目标与评价：相反方向的界迫使等号成立；易错：只套用不等式而未说明为什么必须取等，或漏掉交换后的有序数对。难度：较难。
\newpage
\begingroup\small\setlength{\parskip}{3pt}
\begin{center}{\Large 后两题解析与逐级提示}\end{center}
\textbf{4．证明。（22分）}

十个元素各有“选入”和“不选入”两种选择，因此 $A$ 有 $2^{10}-1=1023$ 个非空子集。每个非空子集的元素和都是正整数，且不超过 $10\times99=990$。

若所有非空子集的和两两不同，则1023个和要占据至多990种数值，这是不可能的。因此存在两个不同的非空子集 $U,V$，其元素和相等。

取 $B=U\setminus V$，$C=V\setminus U$。从两边减去公共部分 $U\cap V$ 的元素和，可知 $B,C$ 的元素和仍相等，且 $B\cap C=\varnothing$。

还必须证明非空。若 $B=\varnothing$，则 $U\subsetneq V$（因为 $U\ne V$），由所有元素为正数，$V$ 的元素和严格大于 $U$ 的元素和，矛盾。同理 $C\ne\varnothing$。证毕。

评分：子集计数6分；和的范围及存在同和子集6分；去掉交集6分；验证两者非空4分。

目标与评价：集合运算、存在性命题与反证；评价重点是“同和”到“不交且非空”的推理是否完整。易错：找到同和子集即结束；忘记正数条件保证差集非空。难度：较难。

逐级提示（学生独立尝试后按需给出）：
\par （1） 一共有多少个非空子集？它们的元素和有多少种可能？
\par （2） 若每个子集的和都不同，数量上会怎样？
\par （3） 两个同和子集可能有公共元素，能否同时删去？为什么不会删空其中一个？

\textbf{5．证明。（24分）}

令 $u=x/(x-1)$，$v=y/(y-1)$，$w=z/(z-1)$。因为 $x,y,z\ne1$，这些分式有意义，且
\[u-1=\frac1{x-1},\quad v-1=\frac1{y-1},\quad w-1=\frac1{z-1}.\]
由 $xyz=1$，得到 $uvw=(u-1)(v-1)(w-1)$。展开、约去 $uvw$，得
\[uv+vw+wu=u+v+w-1.\]
记 $r=u+v+w$，则
\[u^2+v^2+w^2-1=r^2-2(r-1)-1=(r-1)^2\ge0.\]
左边正是待证式左端减去1，故原不等式成立。

核查：$(x,y,z)=(-2,-2,1/4)$ 满足条件且左端等于 $4/9+4/9+1/9=1$，所以常数1可以达到。此例仅用于验证，无须学生另行回答。

评分：有效换元及分母条件6分；转化乘积条件8分；平方恒等式及结论10分。其他正确证明同等给分。

目标与评价：分式结构、等价转化与平方非负；易错：误把实数条件当成正数，直接对有负值的 $u,v,w$ 使用基本不等式。难度：挑战。

逐级提示：
\par （1） 将三个反复出现的分式分别记为 $u,v,w$。
\par （2） 比较 $uvw$ 与 $(u-1)(v-1)(w-1)$。
\par （3） 尝试用 $u+v+w$ 表示 $u^2+v^2+w^2-1$。
\newpage
\endgroup
\begin{center}{\Large 逐题来源、改动与复核记录}\end{center}
\textbf{来源说明（核查日期：2026年9月22日）}

第1、3题：2011年全国高中数学联合竞赛一试A卷第1、3题。核查材料为公开镜像的《试题参考答案及评分标准（A卷）》PDF第1页。页面同时载有题干与解析，已对照页面核查数学条件。
\par\href{https://oss.simiyun.cn/upload/20221115/9f38c37c8c914d33a3e7a04cded2efa7.pdf}{2011年原卷参考答案镜像（点击访问）}

第2题：2014年全国高中数学联赛一试第2题。核查材料为历年试卷汇编PDF第38页，页眉为“2014全国高中数学联赛试题”。该页未标A/B卷，本卷按页面标注，不额外补写卷别。
\par\href{https://7064135.s21i.faiusr.com/61/ABUIABA9GAAgr9Ko6gUo1-Prrgc.pdf}{历年联赛试卷汇编镜像（第38页）}

第4题：1972年第14届国际数学奥林匹克第1题。核查IMO官网英文试题PDF第1页，编号1972/1。
\par\href{https://www.imo-official.org/assets/documents/problems/1972/1972_eng.pdf}{IMO 1972官方试题}

第5题：2008年第49届国际数学奥林匹克第2题（a）。核查IMO官网英文试题PDF第1页，Problem 2(a)。
\par\href{https://www.imo-official.org/assets/documents/problems/2008/2008_eng.pdf}{IMO 2008官方试题}

\textbf{与原题的区别}
\par 1．原填空题改为解答题，补充“说明理由”；条件与数值不变。
\par 2．原题只填写 $M-m$；改为写出 $M,m$ 及差并证明，条件不变。
\par 3．原题询问 $\log_a b$，原答案为 $-1$。为避开尚未学习的对数，保留全部原条件，改问所有有序数对 $(a,b)$；不宣称设问原封不动。
\par 4．中文转述并改写为集合语言；明确“非空”，避免把空集当作平凡答案。
\par 5．只保留（a）的不等式证明；不收录（b）关于无穷多组有理数等号解的证明。不能将本题称为完成了整道IMO第2题。

\textbf{独立数学复核}
\par 1．回代四个三元素子集的和，得到 $-1,3,5,8$。
\par 2．另由 $3/a+a-2\sqrt3=(a-\sqrt3)^2/a\ge0$ 验证下界；最大值、最小值的取等点均在约束内。
\par 3．独立消元得到 $s^2-8p^2=4p(p-1)^2\le0$；由 $p>0$ 必有 $p=1$。回代两组解均满足条件。
\par 4．$1023>990$；正性确保删去交集后两个差集均非空，无需学过排列组合公式。
\par 5．逐项展开验证平方恒等式，并核验等号例子；全程没有假设 $x,y,z$ 为正数。

逻辑教学落点主要在第4题的“存在”、反证及非空性论证；本卷未单独考查充分必要条件分类或含量词命题的否定。若用于单元测验，不能代替这些知识点的常规覆盖。
\end{document}
